% To be reformatted in the TAC style (tac.cls) upon acceptance.
\documentclass[11pt]{amsart}

\usepackage{amsmath,amssymb}
\usepackage[hidelinks]{hyperref}

\newtheorem{thm}{Theorem}[section]
\newtheorem{lem}[thm]{Lemma}
\newtheorem{cor}[thm]{Corollary}
\newtheorem{prop}[thm]{Proposition}
\theoremstyle{definition}
\newtheorem{defi}[thm]{Definition}
\newtheorem{rem}[thm]{Remark}
\newtheorem{obs}[thm]{Observation}

\newcommand{\I}{\mathbb{I}}
\newcommand{\DM}[1]{\mathrm{DM}(#1)}
\newcommand{\Dbox}{\square_{\mathrm{DM}}}
\newcommand{\aconv}{\equiv}
\newcommand{\nconv}{\not\equiv}
\newcommand{\refl}{\mathsf{refl}}
\newcommand{\Path}{\mathsf{Path}}
\newcommand{\Loop}{\mathsf{Loop}}
\newcommand{\U}{\mathcal{U}}
\newcommand{\io}{0}
\newcommand{\ii}{1}
\newcommand{\hcomp}{\mathsf{hcomp}}
\newcommand{\transp}{\mathsf{transp}}
\newcommand{\cell}[1]{\mathsf{#1}}
\newcommand{\act}[1]{A_{#1}}

\begin{document}

\title[The strict layer of cubical type theory]{The strict layer of
  generic cells in cubical type theory\\ is a wedge of De Morgan
  spheres}

\author{Ryota Shibusawa}
\address{Daiichi Institute of Technology, Japan}
\email{r-shibusawa@daiichi-koudai.ac.jp}

\subjclass[2020]{18N45, 03B38, 18M99}
\keywords{cubical type theory, cube categories, connections, De Morgan
  algebra, coherence, definitional equality, weak $\omega$-groupoids}

\begin{abstract}
  In cubical type theories with a De Morgan interval, the connection
  and reversal operations give types strictly more definitional
  structure than the identity types of Martin-L\"of type theory: many
  groupoid laws that are only propositional in Book HoTT hold
  judgmentally.  We determine exactly how much.  Over a context of
  generic higher cells, the cubes definable by pure reparametrization
  --- without Kan operations --- form, modulo definitional equality
  and naturally in the cube dimension, a \emph{wedge of De Morgan
  spheres}: one boundary-collapsed representable presheaf
  $\Dbox(-,[n])/\partial$ on the De Morgan cube category per cell.
  The action of the cube category is strictly functorial (identities
  and composition hold definitionally), it collapses precisely the
  morphisms that factor through a face, it is faithful elsewhere, and
  distinct cells never interact.  In particular transposition of a
  $2$-cell's coordinates is not definitional --- the exact
  definitional shadow of the Eckmann--Hilton argument being weak ---
  and a square such as $q(i \wedge \neg i)(j)$, all four of whose
  edges are degenerate, is not definitionally constant: the
  non-Boolean nature of the free De~Morgan algebra is visible in
  every dimension.  All positive and negative instances are
  machine-checked in a self-contained proof-assistant kernel, and by
  the metatheory of a companion paper every statement is a decidable,
  assumption-free fact about the algorithmic equality of a minimal
  fragment.
\end{abstract}

\maketitle

\section{Introduction}
\label{sec:intro}

A basic discovery of homotopy type theory is that the identity types
of Martin-L\"of type theory equip every type with the structure of a
weak $\omega$-groupoid \cite{Lumsdaine,vdBergGarner}.  The structure
is \emph{weak}: associativity, units, inverses and the exchange laws
hold only up to higher cells, and the bookkeeping of those higher
cells --- the coherence problem --- is notoriously difficult.

Cubical type theories \cite{CCHM} change the balance.  Their
interval $\I$ carries the operations of a De Morgan algebra
($\wedge$, $\vee$, $\neg$), and \emph{reparametrizing} a cell along
these operations --- forming $\langle i \rangle\, p(i \wedge j)$ and
its relatives, the \emph{connections} --- is a definitional
operation: many laws that Book HoTT can only prove propositionally
hold judgmentally.  Practitioners know this well
\cite{VezzosiMortbergAbel}; what has been missing is a precise
statement of \emph{how much} of the $\omega$-groupoid structure the
definitional equality supplies, and where exactly the strict layer
ends and the weak layer necessarily begins.

This paper answers that question for the layer generated by
\emph{generic cells}.  Fix base types and points, and let $q$ be a
generic $n$-cell --- a variable of the $n$-fold iterated loop type.
The definable cubes over such a context, using only
reparametrization (no Kan operations), carry an action of the
\emph{De Morgan cube category} $\Dbox$: objects $[m]$, morphisms
$[m] \to [n]$ the $n$-tuples of free De Morgan polynomials in $m$
variables, composition by substitution.  This is the cube category
with connections, reversals, degeneracies and diagonals, the
maximal ``classical'' cube category in the taxonomy of
\cite{GrandisMauri,BuchholtzMorehouse}.  Our main theorem
(Theorem~\ref{thm:main}) states:

\begin{quote}
  Over a context of generic cells $q_1, \dots, q_s$ of arities
  $n_1, \dots, n_s$, the pure-reparametrization cubes modulo
  definitional equality form, naturally in $[m]$, the wedge of
  pointed sets
  \[
    \bigvee_{t=1}^{s} \Dbox([m],[n_t])/\partial ,
  \]
  where $\Dbox(-,[n])/\partial$ collapses to the basepoint exactly
  the morphisms that factor through a face of $[n]$.
\end{quote}

\noindent That is: each generic cell spans a \emph{De Morgan
sphere} --- the boundary-collapsed representable --- the spheres of
distinct cells meet only at the common basepoint (the constant
cube), and nothing else is identified.  The action realizing the
isomorphism is \emph{strictly} functorial: the identity tuple acts
as the identity and composition acts as substitution, both
definitionally (machine-checked).

Three consequences delimit the strict/weak boundary exactly:
\begin{enumerate}
\item \emph{Transposition is not strict}: for a generic $2$-cell
  $q$, the squares $\langle i \rangle \langle j \rangle\,q(i)(j)$
  and $\langle i \rangle \langle j \rangle\,q(j)(i)$ are distinct
  morphisms of $\Dbox$, hence not definitionally equal.  This is the
  precise definitional shadow of the Eckmann--Hilton argument being
  weak: any exchange of coordinates requires composition structure,
  which lives outside the reparametrization layer.
\item \emph{No excluded middle in any dimension}: $i \wedge \neg i$
  is not an endpoint of the free De Morgan algebra, so the square
  $\langle i \rangle \langle j \rangle\, q(i \wedge \neg i)(j)$ ---
  all four of whose \emph{edges} are degenerate --- is not
  definitionally constant.  The free De Morgan algebra's failure of
  Booleanness is definitional in every dimension.
\item \emph{Cells do not interact}: no definitional equation relates
  cubes built from distinct generic cells, beyond the shared
  basepoint.  Whiskering, horizontal and vertical composition, and
  the interchange law therefore live in the weak layer \emph{by
  theorem}, not merely by inspection of the rules.
\end{enumerate}

\subparagraph*{Honest positioning.}  The positive half of the
characterization --- that the De Morgan laws hold definitionally on
reparametrized cells --- is by design of the theory: the interval
\emph{is} a free De Morgan algebra and the kernel decides its
equality by normal forms \cite{CCHM}.  The content of this paper is
the \emph{converse} and its structure: nothing else is identified
(faithfulness), what collapses is exactly face-factoring
(the sphere structure), the collapse point is shared across cells
(the wedge), and the action is strictly functorial.  The
one-dimensional case of Theorem~\ref{thm:cn} --- a single generic
path, $n = 1$ --- appeared as a self-contained appendix of the
companion paper \cite{Companion}; the present paper strictly
generalizes it (arbitrary arities, higher and multiple cells, the
sphere and wedge structure, the strict functoriality of the action)
and is written to be read independently.  Beyond that special case
we know of no prior statement or proof of this characterization.  It does not
follow from published normalization theorems: the normalization
result of Sterling--Angiuli \cite{SterlingAngiuli} concerns
\emph{Cartesian} cubical type theory, whose cube category has no
connections or reversals, and open-term normalization for the
De~Morgan calculus is not in the literature.  Our proofs are
elementary given the right lemmas --- they analyze the
normalization-by-evaluation conversion of a concrete kernel --- and
every statement is, by the metatheory of the companion paper
\cite{Companion}, a decidable and assumption-free fact about a
minimal fragment of the calculus.

\subparagraph*{Related coherence questions.}  The theorem is a
coherence result of ``every diagram commutes iff it commutes in the
free structure'' type \cite{MacLane}: a diagram of
reparametrizations of a generic cell commutes definitionally iff it
commutes in the free De Morgan algebra --- and, complementing it,
no diagram involving composition commutes definitionally unless
trivially.  The necessity of the weak layer is not an artifact of
our kernel: by the no-go theorem of \cite{Companion}, forcing
composition-layer equations (such as constant-tube collapse for
$\hcomp$) into the definitional equality of a cubical theory
derives regularity principles that fail in the standard models.
The present paper thus pins the strict/weak boundary from both
sides: the strict side is exactly a wedge of De Morgan spheres, and
crossing the boundary is not a matter of better engineering but is
obstructed in principle.

\subparagraph*{Outline.}  Section~\ref{sec:cube} fixes the De
Morgan cube category.  Section~\ref{sec:calc} fixes the fragment of
cubical type theory and its algorithmic equality, and states the
standing metatheoretic facts imported from \cite{Companion}.
Section~\ref{sec:action} defines the action and proves its strict
functoriality.  Section~\ref{sec:single} proves the single-cell
characterization ($C_n$), Section~\ref{sec:multi} the multi-cell
wedge theorem, and Section~\ref{sec:main} assembles the main
theorem.  Section~\ref{sec:weak} records the weak complement.
Section~\ref{sec:artifact} describes the machine verification.

\section{The De Morgan cube category}
\label{sec:cube}

\begin{defi}\label{def:dm}
  $\DM{i_1,\dots,i_m}$ is the free De Morgan algebra on generators
  $i_1,\dots,i_m$: the free bounded distributive lattice with an
  involution $\neg$ satisfying the De Morgan laws
  ($\neg(x \wedge y) = \neg x \vee \neg y$, $\neg 0 = 1$), and
  \emph{without} the laws $x \wedge \neg x = 0$,
  $x \vee \neg x = 1$.  Its elements have canonical normal forms:
  antichains of clauses over the literals
  $i_k, \neg i_k$ (disjunctive normal form, reduced under
  absorption)~\cite{CCHM}.  We write $0, 1$ for the bounds and call
  them the \emph{endpoints}.
\end{defi}

\begin{defi}[The De Morgan cube category]\label{def:cube}
  $\Dbox$ has objects $[m]$ for $m \ge 0$ and morphisms
  \[
    \Dbox([m],[n]) \;:=\; \DM{i_1,\dots,i_m}^{\,n},
  \]
  an $n$-tuple $\vec\varphi = (\varphi_1,\dots,\varphi_n)$ regarded
  as a substitution; composition is substitution of tuples into
  variables, and the identity on $[n]$ is the tuple of generators
  $(i_1,\dots,i_n)$.
\end{defi}

$\Dbox$ is the classical cube category with faces, degeneracies,
symmetries, diagonals, connections of both kinds, and reversals ---
in the taxonomy of Buchholtz--Morehouse \cite{BuchholtzMorehouse} it
is the cube category of the ``De Morgan'' variety; the study of
cubical sites with connections goes back to Grandis--Mauri
\cite{GrandisMauri}.  Two classes of morphisms play a role below:

\begin{defi}\label{def:face}
  The \emph{$k$-th face at $\varepsilon \in \{0,1\}$} is
  $\delta_k^\varepsilon = (i_1,\dots,i_{k-1},\varepsilon,
  i_k,\dots,i_{n-1}) : [n-1] \to [n]$.  A morphism
  $\vec\varphi : [m] \to [n]$ is \emph{face-factoring} if it factors
  through some $\delta_k^\varepsilon$; equivalently, iff some
  coordinate $\varphi_k$ is an endpoint \emph{of the free algebra}
  (i.e.\ its normal form is the empty or the full antichain).
\end{defi}

\begin{defi}[De Morgan spheres]\label{def:sphere}
  $\Dbox(-,[n])/\partial$ is the pointed presheaf obtained from the
  representable $\Dbox(-,[n])$ by identifying all face-factoring
  morphisms to a single basepoint.
\end{defi}

The subtlety that drives several results below is that
face-factoring is a normal-form condition, not a pointwise one:
$i \wedge \neg i$ takes the value $0$ at both endpoints of $i$, yet
it is \emph{not} the endpoint $0$ of the free algebra, so
$(i \wedge \neg i,\, j)$ is not face-factoring.

\section{The calculus, generic cells, and imported metatheory}
\label{sec:calc}

\subsection{The fragment}

We work in a CCHM-style cubical type theory \cite{CCHM} with a De
Morgan interval; everything below lives in the \emph{minimal
fragment}: $\Pi$, $\Sigma$, path types $\Path$, and a base type,
with transport and homogeneous composition --- no universes, no
glueing, no higher inductive types.  In fact the terms of this paper
use only path types over a base type: the fragment of
\emph{reparametrization},
\[
  p \;\mapsto\; \langle i_1 \rangle \cdots \langle i_m \rangle\,
  p(\varphi_1)\cdots(\varphi_n),
  \qquad \varphi_k \in \DM{i_1,\dots,i_m},
\]
where $p(\varphi)$ denotes path application and
$\langle i \rangle$ path abstraction.  Kan operations
($\transp$, $\hcomp$) are \emph{excluded} from the layer under
study; they are what the weak layer is made of
(Section~\ref{sec:weak}).

\begin{defi}[Generic cells]\label{def:generic}
  Over a base context $A : \U$, $a : A$, the \emph{$n$-fold loop
  type} is
  $\Loop^1 := \Path\,A\,a\,a$ and
  $\Loop^{k+1} := \Path\,(\Loop^k)\,\refl_k\,\refl_k$, where
  $\refl_k$ is the $k$-fold reflexivity tower.  A \emph{generic
  $n$-cell} is a context variable $q : \Loop^n$.  For
  $\vec\varphi \in \Dbox([m],[n])$ we write
  \[
    \act{q}(\vec\varphi) \;:=\;
    \langle i_1 \rangle \cdots \langle i_m \rangle\,
    q(\varphi_1)\cdots(\varphi_n),
  \]
  an \emph{$m$-cube} over the context.
\end{defi}

Cells with fully degenerate boundary (loops) are the right generic
objects for a \emph{pointed} characterization; see
Remark~\ref{rem:boundary} for the expected unpointed refinement.

\subsection{Algorithmic equality, and what we import}
\label{sec:imported}

Definitional equality is taken to be the algorithmic equality
$\aconv$ of a concrete normalization-by-evaluation kernel: terms are
evaluated to weak-head values (path abstractions become closures;
applications of variables become \emph{neutral spines}), and
conversion compares values type-directedly, opening binders at fresh
generic dimensions and comparing interval arguments of neutral path
applications by the canonical antichain normal form of
Definition~\ref{def:dm}.  The kernel is the one of the companion
paper \cite{Companion}, whose artifact this paper's probes extend;
three imported facts are standing hypotheses:

\begin{enumerate}
\item[(K1)] \emph{Endpoint collapse is normal-form exact.}  Path
  application of a \emph{value} at an interval argument $r$ collapses
  to the stored endpoint iff $r$ is an endpoint \emph{of the free De
  Morgan algebra} (the kernel decides $r = 0,1$ by the antichain
  normal form, not syntactically).
\item[(K2)] \emph{Conversion of neutrals is spine-wise.}  Two
  neutral spines are convertible iff they have the same head
  variable, the same length, convertible stored endpoint
  annotations, and, at each path-application stage, interval
  arguments equal in the free De Morgan algebra.
\item[(K3)] \emph{Decidability without assumptions.}  On the minimal
  fragment, evaluation and conversion terminate on well-typed input
  and the algorithmic equality is decidable; no base-theory
  assumption is involved.  (This is the endpoint-blind termination
  theorem of \cite{Companion}.)
\end{enumerate}

(K1) and (K2) are facts about two specific functions of the kernel,
verified there; (K3) is a theorem of \cite{Companion}.  Every
positive and negative instance claimed below is additionally
machine-checked (Section~\ref{sec:artifact}); by (K3) each such
check is a decision, not a test.

\begin{rem}[Robustness]\label{rem:robust}
  The reparametrization layer involves no transport, so the results
  are insensitive to the transport-layer design choices studied in
  \cite{Companion} (in particular to its evaluation-time constancy
  rule): they hold verbatim for a standard CCHM-style conversion.
  We state them for our kernel because there they are
  \emph{unconditional} decidable facts by (K3).
\end{rem}

\section{The action and its strict functoriality}
\label{sec:action}

Reparametrization defines, for each generic $n$-cell $q$ and each
$m$, a function
\[
  \act{q} : \Dbox([m],[n]) \longrightarrow
  \{\text{$m$-cubes over } \Gamma\}, \qquad
  \vec\varphi \mapsto \act{q}(\vec\varphi).
\]

\begin{lem}[Unit law, definitionally]\label{lem:unit}
  $\act{q}(\mathrm{id}_{[n]}) \aconv q$.
\end{lem}
\begin{proof}
  $\act{q}(\mathrm{id}) = \langle i_1 \rangle \cdots \langle i_n
  \rangle\, q(i_1)\cdots(i_n)$, which is $\aconv q$ by $n$
  applications of the $\eta$-rule for path types (conversion is
  type-directed: it opens both sides at generic dimensions and
  compares the resulting applications, which are literally equal
  spines).  Machine probe: \texttt{sqEtaD}.
\end{proof}

\begin{lem}[Composition law, definitionally]\label{lem:comp}
  For $\vec\varphi : [m] \to [n]$ and $\vec\psi : [l] \to [m]$,
  \[
    \bigl(\act{q}(\vec\varphi)\bigr)
    \text{ reparametrized along } \vec\psi
    \;\aconv\; \act{q}(\vec\varphi \circ \vec\psi),
  \]
  where $\vec\varphi \circ \vec\psi$ is substitution in $\Dbox$.
\end{lem}
\begin{proof}
  Path application of a path abstraction is $\beta$-reduction, which
  the evaluator implements as substitution of the interval argument
  into the body; iterating, the left-hand side evaluates to the
  right-hand side's evaluation.  Machine probe:
  \texttt{sqCloneCompD} ($(\langle u \rangle \langle v \rangle\,
  q(u)(v))(i \vee j)(i \wedge j) \aconv q(i \vee j)(i \wedge j)$).
\end{proof}

Thus $\act{q}$ is \emph{strictly} functorial --- the action of
$\Dbox$ on cubes is by definitional equality, with no coherence
cells needed.  This is the precise sense in which the connection
layer of cubical type theory is ``strict algebra'': the entire free
De Morgan theory acts on cells on the nose.

\section{The single-cell characterization}
\label{sec:single}

Fix a generic $n$-cell $q$.  The following computation is the
engine of everything.

\begin{lem}[Value computation]\label{lem:value}
  Open the $m$ binders of $\act{q}(\vec\varphi)$ at generic
  dimensions.  The body $q(\varphi_1)\cdots(\varphi_n)$ evaluates
  as follows.
  \begin{enumerate}
  \item If some $\varphi_k$ is an endpoint of the free algebra, the
    value is the base point $a$: the first such coordinate collapses
    the spine --- by (K1), either the application of $q$ itself
    (yielding a reflexivity tower, through which the remaining
    applications compute), or the application of an already-neutral
    spine (yielding its stored endpoint annotation, again a tower)
    --- and towers evaluate through the remaining arguments to $a$.
  \item Otherwise the value is the full neutral spine
    $q(\varphi_1)\cdots(\varphi_n)$ with all $n$ interval arguments
    retained.
  \end{enumerate}
  In particular the value is $a$ iff $\vec\varphi$ is
  face-factoring (Definition~\ref{def:face}).
\end{lem}
\begin{proof}
  By (K1), the collapse at each stage happens iff the argument is a
  free-algebra endpoint; the case analysis is exhaustive.  Machine
  probes: \texttt{sq2DegFstD} (collapse at the first coordinate),
  \texttt{sq2DegSndD} (collapse of a neutral spine at the second),
  \texttt{sq2DegCrossD} (all face-factoring squares are equal).
\end{proof}

\begin{thm}[$C_n$: the sphere of a single cell]\label{thm:cn}
  For $\vec\varphi, \vec\psi \in \Dbox([m],[n])$ (with boundaries
  agreeing so that the comparison is well-typed):
  \[
    \act{q}(\vec\varphi) \aconv \act{q}(\vec\psi)
    \iff
    \begin{cases}
      \text{both are face-factoring, or}\\
      \text{neither is, and } \vec\varphi = \vec\psi
      \text{ in } \Dbox([m],[n]).
    \end{cases}
  \]
  Hence $\act{q}$ descends to an injection
  $\Dbox([m],[n])/\partial \hookrightarrow
  \{\text{$m$-cubes}\}/{\aconv}$ of pointed sets, the basepoint
  being the constant cube.
\end{thm}
\begin{proof}
  Conversion opens the $m$ binders at generic dimensions and
  compares the bodies of Lemma~\ref{lem:value}.

  ($\Leftarrow$)  If both tuples are face-factoring, both bodies are
  the same variable $a$.  If neither is and the tuples are equal,
  the spines agree stage-wise: same head, same stored endpoint
  annotations (the same towers), and interval arguments equal in
  the free algebra, which the antichain normal form decides (K2).

  ($\Rightarrow$)  Suppose conversion accepts.  If both bodies are
  $a$, both tuples are face-factoring by Lemma~\ref{lem:value}.  If
  both are spines, (K2) forces stage-wise equality of the interval
  arguments in the free algebra, i.e.\ $\vec\varphi = \vec\psi$.
  A variable cannot be convertible with a spine, so the mixed case
  is impossible.
\end{proof}

\begin{cor}[Transposition is not strict]\label{cor:transpose}
  For a generic $2$-cell,
  $\langle i \rangle \langle j \rangle\, q(j)(i) \nconv
   \langle i \rangle \langle j \rangle\, q(i)(j)$.
\end{cor}
\begin{proof}
  $(j,i) \ne (i,j)$ in $\Dbox([2],[2])$, and neither is
  face-factoring.  Machine probe with a well-typedness control.
\end{proof}

\begin{cor}[No excluded middle in dimension $2$]\label{cor:lem}
  $\langle i \rangle \langle j \rangle\, q(i \wedge \neg i)(j)
  \nconv \langle i \rangle \langle j \rangle\, a$, although all
  four edges of the left-hand square are degenerate.
\end{cor}
\begin{proof}
  $i \wedge \neg i$ is not an endpoint of the free algebra
  (its antichain normal form is the singleton clause
  $\{i, \neg i\}$), so $(i \wedge \neg i, j)$ is not
  face-factoring; apply Theorem~\ref{thm:cn} against the (constant)
  face-factoring class.  Machine probes \texttt{sq2LEMCtrlD} and
  its companion guard.
\end{proof}

\begin{rem}
  Corollary~\ref{cor:lem} is the higher-dimensional face of a
  familiar dimension-$1$ fact ($\langle i \rangle\, p(i \wedge \neg
  i) \nconv \refl$): the free De Morgan algebra is not Boolean, and
  the definitional equality sees this in every dimension --- a
  square all of whose edges are constant need not be definitionally
  constant.
\end{rem}

\section{Multiple cells: the wedge}
\label{sec:multi}

\begin{obs}[Syntactic exhaustiveness]\label{obs:exhaustive}
  Over a context with generic cells $q_1 : \Loop^{n_1}, \dots,
  q_s : \Loop^{n_s}$ (common base point $a$) and no Kan operations,
  every definable body of an $m$-cube evaluates either to $a$ or to
  an application spine of exactly one $q_t$: the syntax has no
  interval case-split, so a single body term cannot mix two cells.
\end{obs}

\begin{thm}[Wedge theorem]\label{thm:wedge}
  For cells $q_t, q_u$ and tuples $\vec\varphi, \vec\psi$:
  \[
    \act{q_t}(\vec\varphi) \aconv \act{q_u}(\vec\psi)
    \iff
    \begin{cases}
      \text{both tuples are face-factoring
        (any $t, u$), or}\\
      t = u,\ \text{neither is, and } \vec\varphi = \vec\psi.
    \end{cases}
  \]
\end{thm}
\begin{proof}
  Lemma~\ref{lem:value} applies to each cell separately.
  Face-factoring bodies are the \emph{same} variable $a$ regardless
  of the cell --- the shared basepoint.  Non-collapsed spines are
  compared by (K2), which separates distinct head variables
  ($t \ne u$), distinct spine lengths ($n_t \ne n_u$), and distinct
  tuples.  Machine probes: \texttt{sqxDegD} (cross-cell collapse:
  $\langle i\rangle\langle j\rangle q(0)(j) \aconv
  \langle i\rangle\langle j\rangle r(i)(1)$), head separation, and
  spine-depth separation (a $1$-cell square built with the nonzero
  formula $(i \wedge \neg i) \wedge (j \wedge \neg j)$ --- all four
  edges degenerate, interior nonzero --- against a $2$-cell
  square).
\end{proof}

\section{The main theorem}
\label{sec:main}

One more lemma extends the characterization from full applications
to all cubes of higher loop types (partial applications).

\begin{lem}[Partial applications]\label{lem:partial}
  An $m$-cube valued in $\Loop^d$ definable in the layer is a
  partial application $\langle \vec i \rangle\,
  q_t(\varphi_1)\cdots(\varphi_{n_t - d})$ (or a tower).
  Type-directed conversion at $\Loop^d$ opens $d$ further binders
  at fresh generic dimensions; since a generic dimension is never a
  free-algebra endpoint, the comparison reduces to
  Theorem~\ref{thm:wedge} for the extended tuples
  $(\varphi_1,\dots,\varphi_{n_t-d}, l_1, \dots, l_d)$, which are
  face-factoring iff the original partial tuples are.
\end{lem}
\begin{proof}
  Conversion at path types is $\eta$-directed: it compares
  applications at a generic dimension.  Iterating $d$ times appends
  $d$ distinct generic dimensions to the spine, none of them an
  endpoint.  Machine probes: the diagonal $[1] \to [2]$
  (\texttt{sqDiagCtrlD}) and the diagonal/anti-diagonal separation.
\end{proof}

\begin{thm}[Main theorem: a wedge of De Morgan spheres]
  \label{thm:main}
  Over a context of generic cells $q_1, \dots, q_s$ of arities
  $n_1, \dots, n_s$ with common base point, the sets of
  pure-reparametrization $m$-cubes modulo algorithmic equality
  assemble, naturally in $[m] \in \Dbox^{\mathrm{op}}$, into the
  wedge of pointed presheaves
  \[
    \{\text{cubes}\}/{\aconv} \;\cong\;
    \bigvee_{t=1}^{s} \Dbox(-,[n_t])/\partial .
  \]
  The isomorphism is induced by the actions $\act{q_t}$; naturality
  is Lemmas~\ref{lem:unit} and~\ref{lem:comp} (the actions are
  strictly functorial); injectivity is Theorems~\ref{thm:cn}
  and~\ref{thm:wedge}; surjectivity is
  Observation~\ref{obs:exhaustive} with Lemma~\ref{lem:partial}.
  All of this is decidable and assumption-free by (K3).
\end{thm}

\begin{rem}[Boundaries; the unpointed refinement]
  \label{rem:boundary}
  Our generic cells are loops: their boundaries are reflexivity
  towers, and it is exactly this degeneracy that produces the
  $/\partial$ collapse (Lemma~\ref{lem:value} computes the collapse
  \emph{through} the towers).  For a generic cell with \emph{generic
  boundary} --- a context supplying fresh lower-dimensional cells
  for each face --- the same analysis should yield the full
  representable $\Dbox(-,[n])$ with no collapse, and gluing the two
  computations along boundary inclusion should give a cofibration
  of presheaves.  We leave this unpointed refinement, and the
  general position of cell contexts among $\Dbox$-computads, to
  future work.
\end{rem}

\section{The weak complement}
\label{sec:weak}

Theorem~\ref{thm:main} is exhaustive for the layer without Kan
operations; everything else about the $\omega$-groupoid structure
of types --- composition, inverses' cancellation, units of
composition, whiskering, interchange, Eckmann--Hilton --- involves
$\hcomp$ and is \emph{weak}.  Two points deserve emphasis.

First, weakness is now a \emph{theorem} rather than an observation:
by Theorem~\ref{thm:wedge}, no definitional equation relates cubes
built from distinct cells, so no candidate ``definitional
composition'' of generic cells exists in the calculus at all; and by
Corollary~\ref{cor:transpose}, even the one-cell exchange of
coordinates fails definitionally.

Second, the boundary is not an engineering accident.  The
companion paper \cite{Companion} proves a no-go theorem: extending
the definitional equality of a cubical theory with
composition-layer collapse (e.g.\ equating a constant-tube
$\hcomp$ with its base) derives regularity principles of the kind
that fail in the standard cubical-sets models, and the failure
arises already at path types.  Combined with
Theorem~\ref{thm:main}, the picture is complete: \emph{the strict
layer of generic cells is exactly a wedge of De Morgan spheres, and
it cannot be enlarged}.

For calibration, the machine-checked strictness catalogue of
\cite{Companion} sits exactly on this boundary: on the strict side,
all connection reparametrizations, $\eta$-laws, and
transport-related equations up to and including a definitional
$\mathsf{J}$ on $\refl$; on the weak side, all $\hcomp$-generated
laws, beginning with $\mathsf{trans}\,\refl\,\refl \nconv \refl$.

\section{Machine verification}
\label{sec:artifact}

All positive and negative instances quoted in this paper are
elaboration-time checks in a self-contained cubical kernel written
in Lean~4 (no dependence on Lean's definitional equality), the
artifact of \cite{Companion}: repository
\url{https://github.com/r-shibusawa/cubical-strict-j}.
The probes of this paper (\texttt{LibSquares.lean}, together with
the relevant rows of \texttt{LibStrictness.lean} and the supporting
proof document \texttt{docs/ReparamCoherence2.md}) are included from
release \texttt{v1.2.0}, archived at DOI
\href{https://doi.org/10.5281/zenodo.21638543}%
{\texttt{10.5281/zenodo.21638543}}; the companion paper's own snapshot
is release \texttt{v1.1.0}, DOI
\href{https://doi.org/10.5281/zenodo.21637898}%
{\texttt{10.5281/zenodo.21637898}}.
A successful build replays every check.  The probes are decision
procedures by (K3): each catalogue row is a decided fact about the
algorithmic equality, not a test that might be flaky.

The division of labor is the usual one for this artifact: the
kernel supplies machine-checked \emph{instances} (both
identifications and separations); the theorems of this paper are
pen-and-paper results about the conversion algorithm, whose
load-bearing clauses (K1), (K2) are small, inspectable functions.

\section*{Acknowledgements}

The metatheoretic platform (the kernel, the strictness catalogue,
and the decidability theorem for the minimal fragment) was
developed in the companion work \cite{Companion}, whose anonymous
review rounds indirectly shaped the present formulation.

\begin{thebibliography}{10}

\bibitem{CCHM}
C.~Cohen, T.~Coquand, S.~Huber, A.~M\"ortberg.
\newblock Cubical type theory: a constructive interpretation of the
  univalence axiom.
\newblock \emph{TYPES 2015}, LIPIcs 69, 2018.

\bibitem{GrandisMauri}
M.~Grandis, L.~Mauri.
\newblock Cubical sets and their site.
\newblock \emph{Theory and Applications of Categories} 11 (2003),
  185--211.

\bibitem{BuchholtzMorehouse}
U.~Buchholtz, E.~Morehouse.
\newblock Varieties of cubical sets.
\newblock \emph{RAMiCS 2017}, LNCS 10226; arXiv:1701.08189.

\bibitem{Lumsdaine}
P.~L.~Lumsdaine.
\newblock Weak $\omega$-categories from intensional type theory.
\newblock \emph{Logical Methods in Computer Science} 6(3), 2010.

\bibitem{vdBergGarner}
B.~van den Berg, R.~Garner.
\newblock Types are weak $\omega$-groupoids.
\newblock \emph{Proceedings of the London Mathematical Society}
  102(2), 2011, 370--394.

\bibitem{MacLane}
S.~Mac Lane.
\newblock Natural associativity and commutativity.
\newblock \emph{Rice University Studies} 49(4), 1963, 28--46.

\bibitem{VezzosiMortbergAbel}
A.~Vezzosi, A.~M\"ortberg, A.~Abel.
\newblock Cubical Agda: a dependently typed programming language
  with univalence and higher inductive types.
\newblock \emph{Journal of Functional Programming} 31, 2021.

\bibitem{SterlingAngiuli}
J.~Sterling, C.~Angiuli.
\newblock Normalization for cubical type theory.
\newblock \emph{LICS 2021}.

\bibitem{Companion}
R.~Shibusawa.
\newblock Evaluation-time constancy inference for transport in
  cubical type theory.
\newblock Preprint, 2026.  Artifact:
  \url{https://github.com/r-shibusawa/cubical-strict-j}, archived at
  DOI \href{https://doi.org/10.5281/zenodo.21637898}%
  {10.5281/zenodo.21637898}.

\end{thebibliography}

\end{document}
